\documentclass[11pt]{article}
\usepackage[T1]{fontenc}
\usepackage{lmodern,microtype}
\usepackage[margin=1.06in]{geometry}
\usepackage{amsmath,amssymb,amsthm,mathtools}
\usepackage[hidelinks]{hyperref}
\hypersetup{pdftitle={Boundedness of higher-dimensional twisted paraproducts}}
\newtheorem{theorem}{Theorem}[section]
\newtheorem{proposition}[theorem]{Proposition}
\newtheorem{lemma}[theorem]{Lemma}
\newtheorem{corollary}[theorem]{Corollary}
\newcommand{\R}{\mathbb R}
\newcommand{\Z}{\mathbb Z}
\newcommand{\E}{\mathbb E}
\newcommand{\ind}{\mathbf 1}
\newcommand{\dd}{\,\mathrm d}
\newcommand{\avg}[2]{\langle #1\rangle_{#2}}
\newcommand{\anorm}[1]{\lVert #1\rVert}
\newcommand{\A}{\mathcal A}
\newcommand{\G}{\mathcal G}
\newcommand{\B}{\mathcal B}
\newcommand{\D}{\mathcal D}
\DeclareMathOperator{\supp}{supp}
\title{Boundedness of higher-dimensional twisted paraproducts}
\author{}
\date{November 23, 2026}
\begin{document}
\maketitle

\begin{abstract}
We prove strong estimates for the higher-dimensional twisted paraproduct
throughout the finite-input Banach range.
For example, for arbitrary smooth angular kernels of mean zero the
operator maps \((L^{2d}(\mathbb R^d))^d\) into \(L^2(\mathbb R^d)\),
uniformly in both hard radial truncation scales. The proof combines
stopping times for one-coordinate martingales with a recursion on
reflected coordinates. A residual after stopping has small amplitude.
At the step where two copies of this residual occur, fiberwise density
bounds for the tester control their joint contribution.
\end{abstract}

\section{The higher-dimensional twisted paraproduct}
Let \(d\ge2\) and write \([d]=\{1,\ldots,d\}\). For a multiplier
\(m\) on \(\R^d\setminus\{0\}\) and Schwartz functions
\(f_i:\R^d\to\mathbb C\) define
\begin{equation}\label{eq:mult}
 T_m(\vec f)(x)=\int_{(\R^d)^d}
    m(\eta)\prod_{i=1}^d\widehat f_i(\xi^i)
       e^{2\pi i x\cdot(\xi^1+\cdots+\xi^d)}
                  \dd\xi^1\cdots\dd\xi^d,\qquad \eta_i=\xi_i^i .
\end{equation}
Thus \(m\) sees just one frequency coordinate from each input. We use the Mihlin seminorm
\[
  [m]_s=\max_{|\alpha|\le s}\sup_{\eta\ne0}
                        |\eta|^{|\alpha|}|\partial^\alpha m(\eta)|.
\]
Let \(e_1,\ldots,e_d\) denote the standard coordinate vectors.
\begin{theorem}\label{thm:main}
Let \(1<p_i<\infty\), \(1/p=\sum_{i=1}^d1/p_i<1\). There exists
a finite integer \(s=s(d,\vec p)\) such that
\begin{equation}\label{eq:bound}
       \|T_m(\vec f)\|_{p}\le C_{d,\vec p}[m]_s
                        \prod_{i=1}^d\|f_i\|_{p_i}
\end{equation}
for every smooth Mihlin symbol \(m\) and all Schwartz inputs
\(f_1,\ldots,f_d\), with \(T_m\) the untruncated operator in \eqref{eq:mult}.
If \(K(y)=|y|^{-d}\Omega(y/|y|)\) with \(\Omega\in C^\infty(S^{d-1})\)
and \(\int_{S^{d-1}}\Omega=0\), the hard radial truncations
\begin{equation}\label{eq:hard}
  T^K_{\epsilon,R}(\vec f)(x)=\int_{\epsilon<|y|<R}
       K(y)\prod_{i=1}^d f_i(x-y_i e_i)\dd y
\end{equation}
satisfy
\[
    \sup_{0<\epsilon<R<\infty}\|T^K_{\epsilon,R}(\vec f)\|_p
            \le C_{d,\vec p,\Omega}\prod_{i=1}^d\|f_i\|_{p_i}.
\]
The truncations converge in \(L^p\) as \(\epsilon\downarrow0\), \(R\uparrow\infty\).
The uniform estimate for hard truncations also holds for the
off-origin kernel of any smooth Mihlin multiplier, with the constant
controlled by finitely many Mihlin seminorms.
\end{theorem}
\paragraph{Previous estimates.}
The planar operator was Case~6 in Demeter and Thiele's classification
of two-dimensional bilinear Hilbert transforms \cite[Introduction]{DemeterThiele}.
Bernicot showed how to lower the input exponents once one estimate is
known \cite[Theorem~2]{Bernicot}. Kova\v{c} proved the first bounds,
covering \(1/2<p<2\) for finite \(p_1,p_2>1\)
\cite[Theorem~1 and (1.9)]{Kovac}. Lin and Slav\'ikov\'a later allowed
mean-zero \(L^q\) angular kernels, \(q>1\), for \(1\le p<2\), and
identified the smooth-kernel estimate at \(p\ge2\) as open
\cite[Theorem~1.1 and (1.4)]{LS}. Durcik and Roos formulated
\eqref{eq:mult} in all dimensions in connection with averages of simplex
Hilbert transforms. They obtained \(1<p<4/3\), \(1<p_i<4\) in dimension
three and recorded that no tuple was then known for \(d\ge4\)
\cite[(1.10) and Section~3]{DR}. Anisotropic analogues were proved in
dimensions two and three \cite[Theorem~2]{CDR}, \cite[Theorem~4]{DFR}.

\paragraph{Where the planar proof changes.}
In Kova\v{c}'s dyadic form \(\sum_k\int H(E_k^1 F)(\Delta_k^2G)\),
squaring a vertical Haar coefficient introduces two independent vertical
samples. The square is a child-minus-parent difference, so the resulting
box averages telescope. His tree estimate measures the tester \(H\) by
its box norm on each leaf \cite[Proposition~4]{Kovac}. On the unit
square take \(H=\ind_{[0,\rho]\times[0,1]}\). Its box norm is
\(\rho^{1/2}\), since two independent horizontal samples must enter
the strip. The estimate \(L^6\times L^6\to L^3\) in
Theorem~\ref{thm:main} calls for the smaller \(\|H\|_{3/2}=\rho^{2/3}\).

We keep the Haar-square identity. Before using it, we split each low
coordinate average into a value recorded at successive oscillations
of size \(\delta\) and an error bounded by \(\delta\).
For the term with all low values recorded, summation by parts pairs the
variation of one path with the amplitude of another. Its cost grows
only polynomially in \(1+\log(1/\delta)\), uniformly in the number of
spatial scales (Lemma~\ref{lem:all-coarse}).

Every other term contains an error. Each new reflection doubles the
copies of every remaining error. A drift term or a pair of Haar
differences can use up two small factors, so
four copies leave a gain. At the first square there are only two.
Here let \(b\) be the coordinate carrying the original difference,
\(a_i\) the input support measures and \(h\le\min_i a_i\) the tester
support measure. After pruning, the tester has density \(O(h/a_i)\)
on each active \(i\)-fiber \eqref{eq:fiber}. For an error in coordinate
\(i\), one tester copy can be averaged in a fresh \(b\)-variable and
the other in the \(i\)-variable on which the retained coefficient is
constant. Their densities turn its mass \(O(a_i)\) into
\(O(a_i(h/a_b)(h/a_i))=O(h^2/a_b)\). The initial Cauchy--Schwarz
inequality \eqref{eq:first} takes the square root and multiplies by
\(\sqrt{a_b}\), giving \(O(h)\). Set \(R=\max_i a_i/h\). The later
square roots leave an error \(O(hR\delta^{2^{-d}})\) in \eqref{eq:resbound}.
Taking \(\delta\asymp R^{-2^d}\) makes this \(O(h)\). The recorded term
then costs \(O_{\varepsilon}(hR^\varepsilon)\) for every \(\varepsilon>0\)
by Lemma~\ref{lem:all-coarse}, enough for the required
\(h\prod_i(a_i/h)^{1/p_i}\).
Fiberwise decomposition and interpolation give the strong estimates,
and Section~\ref{sec:continuum} passes from dyadic averages to multipliers.


\section{A dyadic estimate and tester densities}\label{sec:prune}
Let \(E_k^i\) be dyadic conditional expectation at interval length
\(2^{-k}\) in coordinate \(i\) alone. Write
\(\Delta_k^i=E^i_{k+1}-E^i_k\). For \(u<v\), \(b\in[d]\), and
\(|c_k|\le1\), consider
\begin{equation}\label{eq:dyadic}
 \Lambda_b(\vec F)=\sum_{k=u}^{v-1}c_k
     \int F_0\,\Delta_k^b F_b\prod_{i\ne b,\ i\ge1}E_k^i F_i .
\end{equation}
\begin{proposition}\label{prop:dyadic}
If \(0<t_i<1\) for \(0\le i\le d\) and \(\sum_{i=0}^dt_i=1\), then
\[
       |\Lambda_b(\vec F)|\lesssim_{d,\vec t}
                   \prod_{i=0}^d\|F_i\|_{1/t_i}.
\]
The bound is uniform in the coefficients and truncations.
\end{proposition}
We first prove the restricted estimate for \(0\le F_i\le\ind_{A_i}\),
with \(a_i=|A_i|<\infty\). Multilinearity then gives the restricted
estimate for complex functions. Null supports are immediate. Write
\(H=F_0,\ h=a_0,\ W=\prod_{i=0}^d a_i^{t_i}\). For fixed endpoints
and \(b\) let \(C\) be the best restricted constant on nonnegative
inputs, uniformly in \((c_k)\). H\"older and contraction give
\(C\le2(v-u)\).

First suppose
\begin{equation}\label{eq:smalltest}
             h\le \min_{1\le i\le d}a_i ,
\end{equation}
and put \(R=\max_{1\le i\le d}a_i/h\).
Let \(M_i\) be the dyadic maximal operator in coordinate \(i\).
Fix \(K>1\), and set
\[
   U_i=\{M_i\ind_{A_0}>K h/a_i\},\qquad G_i=F_i\ind_{U_i^c}.
\]
The fiberwise weak \((1,1)\) estimate gives \(|U_i|\le a_i/K\).
Replacing \(F_i\) by \(G_i\) in telescoping order costs at most
\(d K^{-\min_{i\ge1}t_i}CW\). For a level-\(k\) active cube \(Q\)
define \(H_Q(x)\) as \(H(x)\) times the indicator that for every
\(s\in[d]\) and every dyadic ancestor \(Q'\supseteq Q\)
(including \(Q\)) one has
\[
       \frac1{|I_s(Q')|}\int_{I_s(Q')}\ind_{A_0}(y_s,x_{\hat s})
                                    \dd y_s\le K h/a_s .
\]
Here \(I_s(Q')\) is the \(s\)-interval of \(Q'\) and \(x_{\hat s}\)
denotes the remaining coordinates. On the level-\(v\) cubes set the
new mask to zero. These masks decrease down chains and satisfy
\begin{equation}\label{eq:drop}
   \sum_{k=u}^{v-1}\sum_{Q'\text{ at }k+1}
       \int_{Q'}(H_{\operatorname{par}(Q')}-H_{Q'})\le h .
\end{equation}
At each active \(Q\), for each fixed \(x_{\hat s}\), the mask is zero
along \(I_s(Q)\) if the test for \(s,Q\) fails. Otherwise \(H_Q\le
\ind_{A_0}\) and the density bound holds there. In particular,
for \(I\subseteq I_s(Q)\) of length at least \(2^{-(d+1)}|I_s(Q)|\),
\begin{equation}\label{eq:fiber}
          \frac1{|I|}\int_I H_Q(y_s,x_{\hat s})\dd y_s
                                      \le 2^{d+1}K h/a_s .
\end{equation}

A failed ancestor test in coordinate \(s\) puts the entire fiber
\((I_s(Q),x_{\hat s})\) in \(U_s\), so \(G_s\) is zero along that
interval. If \(s=b\), this kills \(\Delta_k^bG_b(x)\). Otherwise it
kills \(E_k^sG_s(x)\). Thus \(H\) may be replaced by \(H_Q\) in the
pruned form. We will use this replacement also after decomposing the
low paths. For that use their residuals must vanish whenever the
corresponding average is zero.


\section{Stopping the coordinate martingales}\label{sec:stopping}

Fix integers \(u<v\), a distinguished coordinate \(b\), and
\(0\le G_i\le\ind_{A_i}\) with \(a_i=|A_i|<\infty\). Write
\(0\le H\le\ind_{A_0}\), \(0<h=|A_0|\le\min_{i\in[d]}a_i\), and
\(R=\max_{i\in[d]}a_i/h\). We can take the coefficients \(c_k\) real
by splitting their real and imaginary parts. All martingales below are
finite. We run the low martingales \(X^i_k=E^i_kG_i\), \(i\ne b\),
from \(u\) to \(T=v+d+1\), so their values are available on every
staggered interval used in the box recursion. For the high martingale
put \(Y_u=0\) and \(Y_{k+1}-Y_k=c_k\Delta^b_kG_b\) for \(u\le k<v\).
For every \(1<s<\infty\) the scalar martingale transform inequality \cite{Burkholder} gives
\(\anorm{Y_v}_s\lesssim_s a_b^{1/s}\).

For a path \(M=(M_k)_{u\le k\le t}\) write
\(\operatorname{TV}(M)=\sum_{k=u}^{t-1}|M_{k+1}-M_k|\).
Given \(0<\alpha\le1\), set \(Z^\alpha_u(M)=M_u\).
Keep the last recorded value until the first subsequent time at which
the distance from it is at least \(\alpha\), record the new value, and
repeat. For the nonnegative martingales \(X^i\) also record a zero
whenever the last recorded value is nonzero. Almost surely a
nonnegative martingale that reaches zero stays there on the remaining
finite time interval. In particular,
\begin{equation}\label{eq:stop-tracking}
 e^\alpha_k(M)=M_k-Z^\alpha_k(M),\qquad
 |e^\alpha_k(M)|\le\alpha,\qquad
 e^\alpha_k(X^i)=0\ \text{wherever }X^i_k=0 .
\end{equation}
The recorded path and the error are adapted to the same one-coordinate
filtration as \(M\).

\begin{lemma}\label{lem:stop-var}
Let \(Z=Z^\alpha(M)\), where \(M=X^i\) or \(M=Y\), and let \(a=a_i\)
or \(a=a_b\), respectively. Uniformly in the endpoints and in \(\alpha\),
\begin{align}
 \anorm{\alpha\operatorname{TV}(Z)}_q
        &\lesssim_q a^{1/q}\qquad(1\le q<\infty),\label{eq:stop-var}\\
 \int_{\R^d}\sum_k|Z_{k+1}-Z_k|^2&\lesssim a .\label{eq:stop-squares}
\end{align}
For \(M=X^i\), on a generation-\(k\) interval \(I\) in coordinate \(i\)
write the two child error values as
\(e^\alpha_{k+1}=e^\alpha_k+p_k+h_I d_k\), where \(p_k,d_k\) are
constant in \(I\) with the other coordinates fixed. Then
\begin{equation}\label{eq:stop-drift}
 |p_k|\le2\alpha,\quad |d_k|\le\alpha,\quad
 \sum_k\anorm{p_k}_1\lesssim a_i/\alpha,\qquad
 \sum_k\anorm{d_k}_2^2\lesssim a_i .
\end{equation}
\end{lemma}
\begin{proof}
Include the initial time and the terminal time in the sequence of
recording times \(\tau_j\), repeating the terminal time if necessary.
The last, artificially included terminal value need not be assigned to
the recorded path. Put
\[
                V=\sum_j|M_{\tau_{j+1}}-M_{\tau_j}|^2 .
\]
The recording times are bounded stopping times. On each initial dyadic
interval in the active coordinate, with the others fixed,
\((M_{\tau_j})_j\) is a martingale for the stopped filtration. The
scalar martingale square-function inequality at exponent \(2q\) gives
\(\anorm{V}_q\lesssim_q\anorm{M_t}_{2q}^2\). This also holds on
\(\R^d\): raise the local estimate to the \(q\)-th power, sum over
the disjoint initial intervals and integrate the fixed coordinates.
Here \(\anorm{M_t}_{2q}^{2q}\lesssim_q a\) by contraction or by the
martingale transform bound. Each genuine jump of \(Z\) occurs among
the differences defining \(V\), proving \eqref{eq:stop-squares}.

A jump of size \(s\ge\alpha\) contributes \(\alpha s\le s^2\) to
\(\alpha\operatorname{TV}(Z)\). The optional jump to zero of size
less than \(\alpha\) occurs at most once. If a regular jump occurred
earlier, its square covers this additional cost, which is at most
\(\alpha^2\). Otherwise the cost is \(\alpha|M_u|\). Consequently
\(\alpha\operatorname{TV}(Z)\le2V+\alpha|M_u|\), which proves
\eqref{eq:stop-var}.

Finally \(p_k=-E_k^i(Z_{k+1}-Z_k)\). Its \(L^1\) estimate follows
from \eqref{eq:stop-var}. Since
\[
 |d_k|^2\le E_k^i|e^\alpha_{k+1}-e^\alpha_k|^2
 \le 2E_k^i\bigl(|X^i_{k+1}-X^i_k|^2+|Z_{k+1}-Z_k|^2\bigr),
\]
martingale orthogonality and \eqref{eq:stop-squares} give the last
bound in \eqref{eq:stop-drift}. The pointwise estimates follow from
\eqref{eq:stop-tracking} and the two-child representation.
\end{proof}

\begin{lemma}\label{lem:all-coarse}
For any integer \(N\ge0\), put \(\delta=2^{-N}\). For \(1<q<\infty\),
\[
 \left|\int_{\R^d}H\sum_{k=u}^{v-1}
       \left(\prod_{i\ne b} Z^\delta_k(X^i)\right)(Y_{k+1}-Y_k)\right|
 \lesssim_{d,q} (N+1)^d h R^{1/q}.
\]
\end{lemma}
\begin{proof}
For each of the \(d\) paths use the layers \(M^{[0]}=Z^1(M)\) and
\(M^{[n]}=Z^{2^{-n}}(M)-Z^{2^{1-n}}(M)\) for \(1\le n\le N\).
They sum to \(Z^\delta(M)\). For all positive levels,
\(\sup_k|M_k^{[n]}|\le3\,2^{-n}\) by \eqref{eq:stop-tracking}.
For a low path this bound, with the same constant, holds also at
level zero, since \(0\le X^i\le1\). The high path at level zero
may be unbounded. For either type of path, \eqref{eq:stop-var} gives
\begin{equation}\label{eq:stop-layer}
    \anorm{\operatorname{TV}(M^{[n]})}_q\lesssim_q 2^n a^{1/q},
    \qquad 0\le n\le N.
\end{equation}
In particular every occurrence of a variation in the following
integrals obeys
\[
 \int_{\R^d} H\,\operatorname{TV}(M^{[n]})
 \le h^{1-1/q}\anorm{\operatorname{TV}(M^{[n]})}_q
 \lesssim_q 2^n h R^{1/q}.
\]

First take \(Z^\delta(Y)\) in place of \(Y\). In its layer expansion
fix low levels \(n_i\), let \(m\) be the high level, and write
\(L_k=\prod_{i\ne b}(X^i)^{[n_i]}_k\). If \(m\le\max_{i\ne b}n_i\),
bound the sum pointwise by \(\sup|L|\,\operatorname{TV}(Y^{[m]})\).
After integration this costs at most a constant times
\(2^{m-\sum_i n_i}hR^{1/q}\le hR^{1/q}\), using only the variation
estimate for \(Y^{[0]}\) when \(m=0\).

If \(m>\max_{i\ne b}n_i\), then \(m\ge1\). For real paths \(L,W\)
on the finite time interval, discrete summation by parts gives
\[
 \left|\sum_{k=u}^{v-1}L_k(W_{k+1}-W_k)\right|
 \le \sup_k|W_k|\bigl(2\sup_k|L_k|+\operatorname{TV}(L)\bigr).
\]
Apply this with \(W=Y^{[m]}\) and use the product bound
\[
\operatorname{TV}(L)\le
\sum_{r\ne b}\operatorname{TV}((X^r)^{[n_r]})
\prod_{i\ne b,r}\sup_k|(X^i)_k^{[n_i]}| .
\]
The endpoint terms cost \(O_d(h)\). For the variation in slot \(r\) the
bound, up to a constant, is
\[
             2^{-m+n_r-\sum_{i\ne b,r}n_i}hR^{1/q}\le hR^{1/q}.
\]
There are \((N+1)^d\) choices of the levels.

For the remaining path \(W=Y-Z^\delta(Y)\) we have \(\sup|W|\le\delta\).
Take \(L=\prod_{i\ne b}Z^\delta(X^i)\) in the same summation-by-parts
bound. All low factors lie between zero and one, so
\(\operatorname{TV}(L)\le\sum_{i\ne b}\operatorname{TV}(Z^\delta(X^i))\).
H\"older's inequality on \(A_0\) and \eqref{eq:stop-var} bound the
result by \(O_{d,q}(hR^{1/q})\).
\end{proof}


\section{An estimate for the residual terms}\label{sec:residual}
Fix \(0<\delta=2^{-N}\le1/2\), and a nonempty subset
\(D\subseteq[d]\setminus\{b\}\). For \(i\ne b\) let the low path \(L_i\)
be \(e_i^\delta=X_i-Z_i^\delta\) if \(i\in D\), and \(X_i\) otherwise.
Here \(X_i=E^i G_i\) is the martingale of the pruned input from
Section~\ref{sec:prune}. Write \(\Lambda(L)\) for the form on the
left in the following estimate. We will prove
\begin{equation}\label{eq:resbound}
 \left|\sum_{k=u}^{v-1}c_k\int H_Q\,\Delta_k^bG_b
                     \prod_{i\ne b}L_{i,k}\right|
 \ \lesssim_{d,K}\ h+hR\,\delta^{\,2^{-d}} .
\end{equation}
Inside each integral \(Q\) denotes the level-\(k\) cube containing the
integration point. Extend all low paths to level \(v+d+1\) to allow
the offsets in the calculation below.

For an interval \(I\) in coordinate \(i\), write \(L_{i,I}\) for its
path value, a function of the other coordinates. Its values on the two
children are \(L_{i,I}+p_{i,I}\pm d_{i,I}\).
Write \(\sigma_i=\delta\) if \(i\in D\), and \(\sigma_i=1\) otherwise.
Then \(|L_{i,I}|,|d_{i,I}|\le\sigma_i\) and
\(|p_{i,I}|\le2\sigma_i\). On a true martingale \(p_{i,I}=0\).
By \eqref{eq:stop-drift} for residuals and by orthogonality for true
martingales we have
\begin{align}
 \sum_k\int |d_{i,k}|^2&\lesssim a_i, &
 \sum_k\int |p_{i,k}|&\lesssim a_i/\delta\quad(i\in D).\label{eq:dp}
\end{align}

\paragraph{Box energies.}
Let \(b\in J\subseteq[d]\) and \(n=2^{|J|}\). On a dyadic rectangle
\(V=\prod I_i\), choose two independent uniform samples \(x_j^0,x_j^1\)
in \(I_j\) for each \(j\in J\) and one \(x_i\) in \(I_i\) otherwise,
all mutually independent. Write \(x^\omega\) for the corner selected
by \(\omega\in\{0,1\}^J\). With the mask of the cube \(Q\) fixed, put
\[
 B_J(V;Q)=\E\prod_\omega
       \left[ H_Q(x^\omega)\prod_{i\notin J} L_{i,I_i}(x^\omega)\right].
\]
Let \(S_i\) average a function of the rectangle over its two
\(i\)-children, holding \(Q\) fixed, and write \(\partial_i=S_i-I\).
If \(j\in J\), condition on the other coordinates. The two arms of the
product are the same real function \(C\) evaluated on independent
samples. Thus \(B_J\ge0\) and \(\partial_jB_J\ge0\), even with signed
residuals. More precisely the conditional value of \(\partial_j B_J\)
is \(\avg{h_{I_j}C}{I_j}^2\), since
\[
  \tfrac12\sum_{I'\text{ child of }I}\avg C{I'}^2-\avg CI^2
     =\avg{h_I C}{I}^2,\qquad
     h_I=\ind_{I_{\rm left}}-\ind_{I_{\rm right}}.
\]
Here \(\avg{\cdot}{I}\) denotes normalized integration on \(I\).

For fixed offsets \(0\le\nu_i\le |J|-1\) let \(\mathcal V_{k,\nu}\)
be the grid of rectangles with lengths \(2^{-k-\nu_i}\), each with
its enclosing level-\(k\) cube \(Q_k(V)\). Define
\[
 U(J)=\max_{\nu,\,j\in J}
    \sum_{k=u}^{v-1}\sum_{V\in\mathcal V_{k,\nu}}
       |V|\,\partial_j B_J(V;Q_k(V)).
\]
All grid sums at a single level are absolutely integrable, by keeping
one central factor. On the initial grid Cauchy--Schwarz and the
orthogonality of the \(b\)-martingale differences give
\begin{equation}\label{eq:first}
            |\Lambda(L)|\le \sqrt{a_b\,U(\{b\})}.
\end{equation}
Indeed the two factors in Cauchy--Schwarz are the Haar coefficients
in the \(b\)-variable of \(G_b\) and of \(H_Q\prod_{i\ne b}L_{i,I_i}\).
The latter square is exactly \(\partial_bB_{\{b\}}\).

\paragraph{One step of the recursion.}
Write \(\gamma=\prod_{i\notin J}\sigma_i\).
The grid sum of \((\prod_i S_i-I)B_J\) is bounded above by
\(n h\,\gamma^n\). To see this, telescope in \(k\), inserting the
new mask at each child. The initial box has a minus sign and the last
new mask is zero. Only central factors change when replacing inherited
by new functions, since each path is already evaluated at the actual
level of its own interval. In each of the \(n\) resulting differences,
use all low amplitudes, discard the other central factors, and sum the
one mask drop using \eqref{eq:drop}. At a fixed level the staggered
child rectangles partition each new cube, so its integrated drop is
counted once for each choice of the changed corner.

Expand \(\prod_i S_i-I\) as a sum of the individual differences with
preceding refinements, placing the marked coordinate \(j\) first.
Differences in the other reflected coordinates are nonnegative.
To bound \(U(J)\) above we can therefore add the absolute contributions
of the unreflected differences to the preceding boundary bound. The
preceding refinements in a term for \(i\notin J\) produce one grid
whose offsets are at most \(|J|\), permitted for \(U(J\cup\{i\})\).

On such a rectangle put \(A_\omega=L_{i,I_i}(x^\omega)\), and denote
the associated drift and half-difference by \(p_\omega,d_\omega\).
They are independent of the base \(i\)-sample. Let \(C(y)\) be the
product of all central factors and all factors of types other than
\(i\), with the base \(i\)-coordinate set to \(y\).
To compute \(\partial_i B_J\), average over the child signs the
product of \(A_\omega+p_\omega\pm d_\omega\), multiplied by the
child average of \(C\), and subtract the parent expression.
The term using only \(A\)'s cancels. In each expanded term the base
average is \(\avg{h_{I_i}^{\,m}C}{I_i}\), where \(m\) is the
number of selected \(d\)'s. The following bounds are for the absolute
grid contributions, using the outer expectations over all remaining
samples.

\begin{enumerate}
\item If a \(p\) occurs, retain one \(p_\alpha\) and bound the
other \(n-1\) factors of that type by \(2\sigma_i\).
If there is no \(p\) and at least two \(d\)'s, use
\(2|d_\alpha d_\beta|\le |d_\alpha|^2+|d_\beta|^2\) and bound
the remaining \(n-2\) factors by \(\sigma_i\).
All other low types supply \(\prod_{\ell\notin J,\ \ell\ne i}\sigma_\ell^n\).
In both cases \eqref{eq:dp} bounds the sum by a constant times
\begin{equation}\label{eq:multiple}
         a_i\sigma_i^{n-2}
                \prod_{\ell\notin J,\ \ell\ne i}\sigma_\ell^n .
\end{equation}
The drift case has \(\sigma_i=\delta\). To apply \eqref{eq:dp}, the
marginal of the retained corner is uniform on its rectangle and
the \(i\)-levels \(k+\nu_i\) are distinct. This observation applies
to each fixed staggered grid.

\item If exactly one \(d_\alpha\) and no \(p\) occurs, apply
Cauchy--Schwarz over the grid and outer samples. Put \(d_\alpha\)
and the \(n-1\) remaining \(A_\omega\)'s in the first factor.
The second square is the grid sum of \(\avg{h_{I_i} C}{I_i}^2\).
Reflecting the new coordinate \(i\) identifies this with a marked
energy for \(J\cup\{i\}\). By \eqref{eq:dp} the resulting bound is
\begin{equation}\label{eq:single}
               O_d\big(\sigma_i^{n-1}\sqrt{a_i U(J\cup\{i\})}\big).
\end{equation}
\end{enumerate}

At \(J=\{b\}\) and \(i\in D\) we need a better version of
\eqref{eq:multiple}. Retain both central factors. Conditional on the
arguments of the retained \(p_\alpha\) or \(d_\alpha^2\) and on the
base \(i\)-sample \(y\), the opposite corner has a fresh \(b\)-sample.
Its central factor has conditional mean at most \(C_dKh/a_b\) by
\eqref{eq:fiber}. The central factor at \(\alpha\) can then be
integrated in \(y\) with mean at most \(C_dKh/a_i\), independently
of the retained coefficient. Absolute values of the child Haar signs
are one. Thus \eqref{eq:multiple} at this step is at most
\(C_{d,K}h^2/a_b\), since \(n=2\). This is the only step at which a
drift or a double difference of a residual could consume all its
small factors.

\paragraph{Closing the recursion.}
Set \(M=hR=\max_i a_i\). Downward induction using the boundary,
\eqref{eq:multiple}, and \eqref{eq:single} first gives
\(U(J)\lesssim_d M\) for every \(J\), including the terminal
\(J=[d]\). For \(|J|\ge2\) and \(D\setminus J\ne\varnothing\) it gives
the sharper bound
\[
        U(J)\lesssim_d M\delta^{\,2^{-(d-|J|)}} .
\]
Indeed the boundary has a factor \(\delta^n\), and every multiple
term has at least \(\delta^{n-2}\), with \(n\ge4\).
A singleton at a residual type gains \(\delta^{n-1}\) by
\eqref{eq:single} and the crude bound. A singleton at an ordinary
type leaves a residual in the next node, so one takes the square
root of its inductive gain. The terminal implication is vacuous
when all types have been reflected.

At \(J=\{b\}\) the boundary and the multiple terms for ordinary
types retain a factor \(\delta^2\). The multiple terms for residual
types were bounded by \(C_{d,K}h^2/a_b\). For a singleton at a
residual type use its factor \(\delta\), and for an ordinary type
use the next node's gain. Hence
\[
 U(\{b\})\lesssim_{d,K}h^2/a_b+M\delta^{\,2^{-(d-1)}}.
\]
Together with \eqref{eq:first} and \(a_b\le M\) this proves
\eqref{eq:resbound}. This argument also covers \(d=2\), where
the only singleton at the first node is a residual type.


\section{Restricted and strong estimates}
\begin{proof}[Proof of Proposition~\ref{prop:dyadic}]
First suppose \(h\le\min_{i\ge1}a_i\). For the pruned inputs set
\[
       N=\max\{1,\lceil 2^d\log_2 R\rceil\},\qquad \delta=2^{-N}.
\]
On the low slots write \(X_i=Z_i^\delta+e_i^\delta\). We use the identity
\[
 \prod_{i\ne b}X_i-\prod_{i\ne b}Z_i^\delta
     =\sum_{\varnothing\ne D\subseteq[d]\setminus\{b\}}
       (-1)^{|D|+1}\prod_{i\in D}e_i^\delta
                       \prod_{i\notin D,\ i\ne b}X_i .
\]
For the term with all \(Z\)'s keep the original tester \(H\) and use
Lemma~\ref{lem:all-coarse}. In each summand on the right insert \(H_Q\).
This is exact by Section~\ref{sec:prune} and the zero rule
\eqref{eq:stop-tracking}. Now \eqref{eq:resbound} costs \(O_{d,K}(h)\)
per summand since \(R\delta^{2^{-d}}\le1\).

Write \(t_*=\min_{i\ge1}t_i\), and choose \(q>2/t_*\).
Since \(N\lesssim_d 1+\log R\),
\[
   (1+N)^d R^{1/q}\lesssim_{d,\vec t} R^{t_*}
        \le \prod_{i=1}^d (a_i/h)^{t_i}=W/h .
\]
Restoring the discarded inputs now gives
\begin{equation}\label{eq:small-bound}
            |\Lambda_b(\vec F)|
                 \le (C_{d,K,\vec t}+\delta_K C)W,\qquad
                   \delta_K=dK^{-t_*}.
\end{equation}


For arbitrary support sizes fix \(L>d\). For each \(a_i<h/L\) put
\(V_i=\{M_i\ind_{A_i}>L a_i/h\}\).
Restricting the tester to their union costs \(C(d/L)^{t_0}W\).
On each \(i\)-fiber replace \(F_i\) on the maximal intervals of
\(V_i\) by its interval average. Every dyadic average and difference
of the resulting input agrees with that of \(F_i\) at points outside
\(V_i\). Almost every point of \(A_i\) belongs to \(V_i\).
The replacement has support of measure at most \(h/L\) and size at
most \(2L a_i/h\), by the maximal intervals' parent densities.
Divide by this size bound and enlarge the supporting set to measure
\(h\). For \(h/L\le a_i<h\) enlarge the support to measure \(h\)
directly. Use the original support of measure \(h\) for the
remaining tester. This places the normalized data under
\eqref{eq:smalltest}. Undoing the normalization increases the
target relative to the original \(W\) by at most \((2L)^d\).
For example, each replaced input costs a ratio
\(2L(a_i/h)^{1-t_i}\le 2L\).
Equation \eqref{eq:small-bound} implies
\[
 C\le (2L)^d C_{d,K,\vec t}+
          \big[(d/L)^{t_0}+(2L)^d\delta_K\big] C .
\]
First choose \(L\), then \(K\), so that the bracket is at most
\(1/2\). This proves the uniform restricted bound. Every tuple with all
\(t_i>0\) belongs to a relative open neighborhood in the H\"older
hyperplane on which these restricted bounds hold. Multilinear restricted-type interpolation
gives the asserted strong estimate \cite{GT}.
\end{proof}

\section{Continuous multipliers and truncations}\label{sec:continuum}
Take \(t_i=1/p_i\) for \(i\in[d]\) and \(t_0=1-1/p\) for the
dual tester. We give the comparison with \eqref{eq:dyadic} to
keep track of uniformity in the scales and the choice of \(b\).
For a one-dimensional profile \(\phi\) put
\(\phi_k(y)=2^k\phi(2^k y)\), and let \(*_i\) denote convolution in
coordinate \(i\). For the exponents of Proposition~\ref{prop:dyadic}
the separated sum
\begin{equation}\label{eq:smoothcone}
   \sum_{k=u}^{v-1} c_k(\psi_k*_b f_b)
                   \prod_{i\ne b}(\phi_{i,k}*_i f_i)
\end{equation}
is bounded in \(L^p\) by the product of the input norms times a
constant controlled by finitely many profile seminorms. Here the
profiles are Schwartz, \(\widehat\psi\) is supported in a fixed
compact annulus, and \(\phi_{i,k}=(\phi_i)_k\).

To prove this, let \(A(y)=(1-|y|)_+\). First replace each
\(\phi_i\) by \((\int\phi_i)A\). The difference profile has integral
zero and a bounded square function on every \(L^q,\ 1<q<\infty\).
Indeed, writing this difference as \(a\), one has
\(\lvert\widehat a(\zeta)\rvert\lesssim
 \min(|\zeta|,|\zeta|^{-1})\) by its finite first moment and
integrable derivative. The kernels \(\sum_k\varepsilon_k a_k\)
have a uniform \(L^2\) multiplier bound and the usual
Calder\'on--Zygmund estimates, uniformly over signs and finite
sums. For the kernel estimates one sums the dilated bounds on
\(a\) and its derivative; the tent is compactly supported and
Lipschitz. The scalar singular integral theorem and Khintchine's
inequality give the square estimate. The same argument applies
to \(\psi\). In the telescoping replacement of the low profiles
bound the two square functions by pointwise Cauchy--Schwarz in
\(k\), and dominate every remaining low factor by a directional
maximal function. H\"older gives the required bound for all errors.

For a finite scale interval, translating a dyadic grid uniformly
over a period of its coarsest scale makes the expected kernel of
averaging on intervals of length \(u_0\) equal to
\(u_0^{-1}A(\cdot/u_0)\). This is the probability that two points
lie in the same interval, divided by its length. Take independent
translates in all coordinates and dilate the \(b\)-grid by a fixed
\(\lambda>0\). The dyadic estimate on these grids follows by
translation and a change of scale in that coordinate.
The expected difference kernel in \(b\) has Fourier multiplier
\(D(2^{-k}\zeta)\), with
\[
     D(z)=\widehat A(\lambda z/2)-\widehat A(\lambda z).
\]
Choose \(\lambda\) small enough that \(D\) has no zeros on
\(\supp\widehat\psi\), using the nonzero quadratic coefficient of
\(\widehat A\) at zero. Let \(V_k\) have one-dimensional symbol
\((\widehat\psi/D)(2^{-k}\zeta)\), smoothly extended by zero.
Apply the dyadic bound using coefficients \(c_k\varepsilon_k\)
and the distinguished input \(\sum_l\varepsilon_l V_l^{\,b}f_b\).
This input has \(L^{p_b}\) norm \(\lesssim\|f_b\|_{p_b}\)
uniformly in the signs, by the scalar multiplier theorem and
finite overlap of its bands. Expectation over independent signs
retains precisely \(l=k\). Averaging the grids now gives
\eqref{eq:smoothcone} with all its low profiles equal to \(A\).
Duality against \(L^{p'}\) justifies the averaging in the norm bound.

Decompose \(m\) by a dyadic annular partition and then by a smooth
partition into regions where \(|\eta_b|\gtrsim|\eta|\).
On rescaling each piece to unit size, expand in a Fourier series
on a fixed cube, with a smooth product cutoff equal to one on
the piece and supported inside the period cube. The cutoff in
coordinate \(b\) may be chosen in a fixed annulus. The coefficient
of mode \(n\in\Z^d\) at any scale is bounded by
\(C_{d,N}[m]_N(1+|n|)^{-N}\) by integration by parts. For this
mode the separated profiles in \eqref{eq:smoothcone} are fixed
up to translations determined by \(n\), and their required seminorms
grow at most polynomially in \(|n|\). Taking \(N\) large enough
allows absolute summation in operator norm. Proposition~\ref{prop:dyadic} applies to every choice of \(b\).
Dominated convergence on the Fourier side for Schwartz inputs
and Fatou's lemma give \eqref{eq:bound}.

For hard truncations let \(\mathcal K=\mathcal F^{-1}m\), with
off-origin kernel \(K\). Smooth Mihlin bounds give
\(|K(y)|+|y||\nabla K(y)|\lesssim |y|^{-d}\), with constants
controlled by finitely many seminorms. Indeed each frequency piece
at radius \(2^k\), differentiated \(\ell=0,1\) times in space, is
bounded by \(C 2^{k(d+\ell)}(1+2^k|y|)^{-d-2}\) by integration
by parts, and the sum is geometric. Take a smooth mollifier \(\rho\)
of mass one supported in the unit ball and write
\(\rho_r(y)=r^{-d}\rho(y/r)\). Uniformly in \(r>0\),
\[
       \|\ind_{\{|y|>r\}}K-\mathcal K*\rho_r\|_1\lesssim 1.
\]
On \(|y|\le2r\) use
\(|\mathcal K*\rho_r|\le\|m\|_\infty\|\widehat{\rho_r}\|_1
\lesssim r^{-d}\) and integrate \(|K|\) on \(r<|y|\le2r\).
On \(|y|>2r\) the mean value theorem bounds the difference by
\(Cr|y|^{-d-1}\). Consequently the hard annulus differs by a uniformly
bounded \(L^1\) kernel from the operator kernel with multiplier
\[
          m(\eta)\big(\widehat\rho(\epsilon\eta)-\widehat\rho(R\eta)\big).
\]
This symbol has uniform Mihlin seminorms. For the \(L^1\) error
Minkowski and H\"older bound its operator norm by its kernel norm.
This proves the uniform estimate for every smooth Mihlin kernel.

For \(K(y)=|y|^{-d}\Omega(y/|y|)\) with smooth mean-zero \(\Omega\),
the principal-value distribution has a smooth Mihlin multiplier.
For completeness, partition \(K\) into smooth radial annular pieces.
The piece at scale \(r\), rescaled to unit size, is smooth and has
integral zero. At \(r|\eta|\le1\) its Fourier transform is
\(O(r|\eta|)\), and its derivatives of order \(n\ge1\) are \(O(r^n)\).
At \(r|\eta|>1\) integration by parts gives
\(O_M(r^n(r|\eta|)^{-M})\) for derivatives of order \(n\ge0\).
Summing proves the stated symbol estimates.

For Schwartz inputs, near zero subtract \(\prod_i f_i(x)\);
the difference is \(O_M(|y|(1+|x|)^{-M})\) for \(|y|\le1\).
For \(|y|\ge1\) the product decays faster than any inverse power
of \(1+|x|+|y|\), since \((x,y)\mapsto(x-y_i e_i)_{i=1}^d\)
is injective for \(d\ge2\). These estimates give norm convergence
on Schwartz data. Multilinearity, density, and the uniform bound
complete the proof of Theorem~\ref{thm:main}.

\begin{thebibliography}{99}
\bibitem{Bernicot}
F.~Bernicot, \emph{Fiber-wise Calder\'on--Zygmund decomposition and
application to a bi-dimensional paraproduct},
Illinois J. Math. \textbf{56} (2012), 415--422.
\href{https://arxiv.org/abs/1011.3614}{arXiv:1011.3614}.
\bibitem{Burkholder}
D.~L.~Burkholder, \emph{Martingale transforms},
Ann. Math. Statist. \textbf{37} (1966), 1494--1504.
\href{https://doi.org/10.1214/aoms/1177699141}{doi:10.1214/aoms/1177699141}.
\bibitem{CDR}
M.~Christ, P.~Durcik, and J.~Roos,
\emph{Trilinear smoothing inequalities and a variant of the triangular
Hilbert transform}, Adv. Math. \textbf{390} (2021), 107863.
\href{https://doi.org/10.1016/j.aim.2021.107863}{doi:10.1016/j.aim.2021.107863}.
\bibitem{DemeterThiele}
C.~Demeter and C.~Thiele,
\emph{On the two-dimensional bilinear Hilbert transform},
Amer. J. Math. \textbf{132} (2010), 201--256.
\href{https://doi.org/10.1353/ajm.0.0101}{doi:10.1353/ajm.0.0101}.
\bibitem{DFR}
P.~Durcik, M.~Fraccaroli, and J.~Roos,
\emph{A curved three-dimensional simplex Hilbert transform},
preprint (2026). \href{https://arxiv.org/abs/2610.24850}{arXiv:2610.24850}.
\bibitem{DR}
P.~Durcik and J.~Roos,
\emph{Averages of simplex Hilbert transforms},
Proc. Amer. Math. Soc. \textbf{149} (2021), 633--647.
\href{https://doi.org/10.1090/proc/15196}{doi:10.1090/proc/15196}.
\bibitem{GT}
L.~Grafakos and T.~Tao,
\emph{Multilinear interpolation between adjoint operators},
J. Funct. Anal. \textbf{199} (2003), 379--385.
\href{https://arxiv.org/abs/math/0111141}{arXiv:math/0111141}.
\bibitem{Kovac}
V.~Kova\v{c}, \emph{Boundedness of the twisted paraproduct},
Rev. Mat. Iberoam. \textbf{28} (2012), 1143--1164.
\href{https://doi.org/10.4171/RMI/707}{doi:10.4171/RMI/707}.
\bibitem{LS}
F.~Y.-H.~Lin and L.~Slav\'ikov\'a,
\emph{Rough averages of triangular Hilbert transforms},
preprint (2026). \href{https://arxiv.org/abs/2607.20206}{arXiv:2607.20206}.
\end{thebibliography}
\end{document}
